\section{Síntesis y ampliación}

\subsection{Puntos clave}

\begin{enumerate}
    \item \textbf{Reasoning Agent} es un paradigma de Agente completamente nuevo que usa el lenguaje natural como representación interna y cuenta con un espacio de razonamiento de dimensión infinita.
    \item \textbf{ReAct} unifica el razonamiento y la acción, que se refuerzan mutuamente: el razonamiento guía la planeación de acciones y el manejo de excepciones, mientras que la acción aporta al razonamiento información externa y cómputo.
    \item La \textbf{memoria de largo plazo} (Reflexion, Voyager, Generative Agents) permite que el Agente acumule experiencia entre tareas.
    \item La \textbf{automatización digital} es una dimensión de aplicación completamente nueva que abren los Agentes LLM.
    \item La simplicidad (simplicity) es la cualidad de investigación más poderosa: simple significa general.
\end{enumerate}

\subsection{Lecturas adicionales}

\begin{itemize}
    \item Yao et al., ``ReAct: Synergizing Reasoning and Acting in Language Models,'' ICLR 2023.
    \item Shinn et al., ``Reflexion: Language Agents with Verbal Reinforcement Learning,'' NeurIPS 2023.
    \item Wang et al., ``Voyager: An Open-Ended Embodied Agent with Large Language Models,'' 2023.
    \item Park et al., ``Generative Agents: Interactive Simulacra of Human Behavior,'' UIST 2023.
    \item Yao et al., ``WebShop: Towards Scalable Real-World Web Interaction with Grounded Language Agents,'' NeurIPS 2022.
    \item Zhou et al., ``WebArena: A Realistic Web Environment for Building Autonomous Agents,'' ICLR 2024.
    \item Yao et al., ``Tau-Bench: A Benchmark for Tool-Agent-User Interaction in Real-World Domains,'' 2024.
\end{itemize}

%% --- Fin del contenido --- %%

\end{document}
